\begin{workedexample}[Worked Example: A 4:1 balun impedance transform]
A balun often also transforms impedance. Matching a $200\ \Omega$ balanced
antenna (e.g.\ a folded dipole) to a $50\ \Omega$ unbalanced feed needs a ratio
$Z_{\text{bal}}/Z_{\text{unbal}} = 200/50 = 4$, i.e.\ a turns ratio of
$\sqrt{4} = 2{:}1$. Beyond impedance, the balun forces equal-and-opposite
currents on the balanced side, suppressing common-mode current on the feedline.
\end{workedexample}

\begin{keytakeaways}
\begin{itemize}[leftmargin=1.4em,itemsep=2pt]
\item A balun couples a balanced load to an unbalanced line and suppresses common-mode current.
\item Impedance ratio $= (\text{turns ratio})^2$: a $2{:}1$ turns balun is $4{:}1$ in impedance.
\item $90^{\circ}$ and $180^{\circ}$ hybrids give fixed phase splits for mixers and balanced amplifiers.
\end{itemize}
\end{keytakeaways}

\begin{exercises}
\item What turns ratio gives a 9:1 impedance balun?
\item What impedance-ratio balun matches a $300\ \Omega$ folded dipole to $75\ \Omega$ coax?
\item What does a $180^{\circ}$ hybrid provide that a Wilkinson does not?
\end{exercises}

\furtherreading{Pozar~\cite{pozar} (ch.~7); Balanis~\cite{balanis}.}
